% ------------------------------------------------------------- Results
% =====================================================================
% Results figures (pgfplots). Data: figdata/*.dat from make_numbers.py
% =====================================================================
\pgfplotsset{
  p4/.style={width=\columnwidth, height=4.6cm, grid=major, grid style={black!12},
    tick label style={font=\scriptsize}, label style={font=\footnotesize},
    legend style={font=\scriptsize, draw=none, fill=white, fill opacity=0.85, text opacity=1},
    every axis plot/.append style={line width=0.9pt}},
}

% ---------------- Fig: detector operating curves ------------------------
\newcommand{\figROC}{%
\begin{figure}[t]
\centering
\begin{tikzpicture}
\begin{axis}[p4, width=0.56\columnwidth, height=4.3cm, xmin=0, xmax=0.3, ymin=0, ymax=1.02,
  xlabel={FPR}, ylabel={TPR}, title={\scriptsize step offset, 5\,m},
  legend to name=rocleg, legend columns=4]
\addplot[cBlue] table[x=fpr,y=tpr]{figdata/roc_step_D1.dat}; \addlegendentry{D1 signature}
\addplot[cOrange, dashed] table[x=fpr,y=tpr]{figdata/roc_step_D2.dat}; \addlegendentry{D2 $\chi^2$}
\addplot[cGreen, dotted] table[x=fpr,y=tpr]{figdata/roc_step_D3.dat}; \addlegendentry{D3 SAVIOR-style}
\addplot[cRed, dashdotted] table[x=fpr,y=tpr]{figdata/roc_step_D4.dat}; \addlegendentry{D4 invariant}
\draw[black!50, densely dashed] (axis cs:0.05,0) -- (axis cs:0.05,1.02);
\end{axis}
\end{tikzpicture}%
\hfill
\begin{tikzpicture}
\begin{axis}[p4, width=0.56\columnwidth, height=4.3cm, xmin=0, xmax=0.3, ymin=0, ymax=1.02,
  xlabel={FPR}, yticklabels={}, title={\scriptsize ramp, 0.5\,m/s}]
\addplot[cBlue] table[x=fpr,y=tpr]{figdata/roc_ramp_D1.dat};
\addplot[cOrange, dashed] table[x=fpr,y=tpr]{figdata/roc_ramp_D2.dat};
\addplot[cGreen, dotted] table[x=fpr,y=tpr]{figdata/roc_ramp_D3.dat};
\addplot[cRed, dashdotted] table[x=fpr,y=tpr]{figdata/roc_ramp_D4.dat};
\draw[black!50, densely dashed] (axis cs:0.05,0) -- (axis cs:0.05,1.02);
\end{axis}
\end{tikzpicture}\\[-1pt]
\pgfplotslegendfromname{rocleg}
\caption{Detector operating curves against two \emph{naive} offsets, with
thresholds swept over \nCalMissions{} benign missions (dashed line: the 5\%
operating point). Only D2 and D4 separate a naive 5\,m step, and only D2
separates the ramp.}
\label{fig:roc}
\end{figure}}

% ---------------- Fig: anchor-interval sweep ----------------------------
\newcommand{\figAnchor}{%
\begin{figure}[t]
\centering
\begin{tikzpicture}
\begin{loglogaxis}[p4, height=4.4cm, xmin=4, xmax=25, ymin=1, ymax=400,  xtick={5,10,20}, xticklabels={5,10,20},
  xlabel={re-anchoring interval $T_a$ (s)}, ylabel={max.\ undetected contribution (m)},
  legend pos=north west, legend cell align=left]
\addplot[cRed, mark=square*, mark size=1.6pt, line width=1.2pt] table[x=Ta,y=maxall]{figdata/anchor_sweep_open.dat}; \addlegendentry{final interval unchecked}
\addplot[cBlue, mark=*, mark size=1.6pt, line width=1.2pt] table[x=Ta,y=maxall]{figdata/anchor_sweep.dat}; \addlegendentry{every interval closed (C2$'$)}

\end{loglogaxis}
\end{tikzpicture}
\caption{Stealth budget against the full stack (5\% FPR) versus the
re-anchoring interval, pooled over 30 attackers (15 tuned per $T_a$ and the 15
optimised against the full stack) and maximised over the five disguise
families, on the same missions at every $T_a$. With every interval closed, the
largest undetected contribution grows from \nAnchorMaxFive{} to
\nAnchorMaxTwenty\,m (\nAnchorMultTwenty$\times$). With the final segment
before mission end unchecked, it grows \nOpenMultTwenty$\times$ to
\nOpenMaxTwenty\,m.}
\label{fig:anchor}
\end{figure}}

% ---------------- Fig: dependence gap -----------------------------------
\newcommand{\figGap}{%
\begin{figure}[t]
\centering
\begin{tikzpicture}
\begin{semilogyaxis}[p4, height=4.8cm, xmin=2.5, xmax=8.5, ymin=1e-4, ymax=1.2,
  xlabel={budget available to benign sources $s$ (m)}, ylabel={$\Pr[\sum_i X_i>s]$},
  legend style={at={(0.5,-0.28)}, anchor=north, legend columns=2}]
\addplot[cPurple, line width=1.2pt] table[x=s,y=union]{figdata/gap.dat}; \addlegendentry{union bound}
\addplot[cPurple, dashed, line width=1.2pt] table[x=s,y=ra]{figdata/gap.dat}; \addlegendentry{RA (worst coupling)}
\addplot[cBlue, dashdotted] table[x=s,y=indep]{figdata/gap.dat}; \addlegendentry{indep., upper envelope}
\addplot[cGreen, densely dotted, line width=1.1pt] table[x=s,y=indepmix]{figdata/gap.dat}; \addlegendentry{indep., matched wind}
\addplot[cOrange, mark=x, mark size=1.6pt, only marks] table[x=s,y=sumsame]{figdata/gap.dat}; \addlegendentry{sum of suprema, same mission}
\addplot[cGray, mark=*, mark size=1pt, only marks] table[x=s,y=joint]{figdata/gap.dat}; \addlegendentry{joint simulation}
\end{semilogyaxis}
\end{tikzpicture}
\caption{Where the looseness lives (tail of the benign budget). Union and RA
nearly coincide. The steps down from RA are: to independent sums at the upper
wind envelope (dependence gap); to independent sums at the wind actually
flown (envelope gap); to the sum of per-source suprema on the same missions,
which confirms that the simulated sources are independent; and to the joint
simulation (temporal gap: sources seldom peak at the same instant).}
\label{fig:gap}
\end{figure}}

% ---------------- Fig: floor vs MCR --------------------------------------
\newcommand{\figFloor}{%
\begin{figure}[t]
\centering
\begin{tikzpicture}
\begin{axis}[p4, height=4.6cm, xmin=4, xmax=25, ymin=0, ymax=1.03,
  xlabel={corridor margin $m$ (m)}, ylabel={$1-\Pr[\mathrm{fail}\wedge U]$},
  legend style={at={(0.5,-0.28)}, anchor=north, legend columns=2}]
\addplot[cPurple, line width=1.3pt] table[x=m,y=ra]{figdata/floor_curve.dat}; \addlegendentry{$1-\pbar$, RA, $\theta_{\mathrm{net}}=0$}
\addplot[cPurple, dashed] table[x=m,y=ranet]{figdata/floor_curve.dat}; \addlegendentry{$1-\pbar$, RA, $\theta_{\mathrm{net}}=0.25$\,s}
\addplot[cBlue, dashdotted] table[x=m,y=indep]{figdata/floor_curve.dat}; \addlegendentry{independence (not certified)}
\addplot[cRed, mark=*, mark size=2pt, only marks] table[x=m,y=obs]{figdata/floor_points.dat}; \addlegendentry{measured $1-\Pr[\mathrm{fail}\wedge U]$}
\end{axis}
\end{tikzpicture}
\caption{Certified bound versus measurement at 5\% FPR. The measured
points are the worst case over five held-out attackers and three wind
conditions. The bound stays below every measurement. Including the network
term shifts it right by $e^{LT_c}\dnet=\nNetTubeQuarter$\,m.}
\label{fig:floor}
\end{figure}}

The pilot comprises \nMissionsTotal{} simulated missions and ran in
\nWallMin{} minutes on two CPU cores. Every number below is generated from the
released result files.

\subsection{Calibration: who sets the thresholds?}
At the 5\% operating point, the per-detector FPR stays between \nFprModMin{}
and \nFprModMax{} under moderate wind, and the anchor gate alone contributes
\nFprAnchor. One finding shapes everything that follows: \emph{wind and faults
alone trip no detector} (FPR \nFprWindFaultMax{} for every detector when either
is the only benign source), whereas GNSS error alone produces up to
\nFprGnssOnlyMax{}. The thresholds are therefore set by the tails of the GNSS
error, and the tolerance they grant is the space an attacker exploits. Severe
turbulence, outside the declared envelope, raises the FPR of the SAVIOR-style
detector D3, which models no wind, to \nFprSevereDThree{}; the others stay at or
below \nFprSevereOtherMax{}.

\Cref{fig:roc} shows the operating curves against naive offsets. At 5\% FPR the
$\chi^2$ gate detects a 5\,m step in \nTprStepDTwo{} of missions and the
invariant D4 in \nTprStepDFour{}, but the signature monitor and D3 detect it in
only \nTprStepDOne{} and \nTprStepDThree{}, which is close to chance. A 0.5\,m/s
ramp is caught by D2 (\nTprRampDTwo{}) and essentially by no other detector
($\le$\nTprRampDThree{}).

\figROC

\subsection{RQ1: evasion is universal; the damage is not}
\paragraph{Cross-evaluation.}
We evaluated each of the \nPooledAttackers{} optimised attackers against every
stack. In \nOwnOptFail{} of \nCells{} (stack, family) cells, the CEM run that
targeted that stack found no parameterisation evading in at least 5\% of
missions, although an attacker optimised for a different stack does. Reading detector strength off the
optimiser's own target would therefore have been wrong, so
\cref{tab:disguise} reports the best attacker of each family, wherever it was
optimised.

\paragraph{With the anchor gate, the anchor sets the budget.}
With condition C2$'$ in force and its gate armed (\cref{tab:disguise}, top),
every one of the five families evades every detector, and the full stack, in
at least \nCrossEvMin{} of missions; for reference, benign missions raise no
alarm in about \nBenignNoAlarm{} under the full stack. Each family has
parameterisations that offset GNSS by about 1--2\,m, and every detector must
tolerate offsets of that size to keep its FPR at 5\%. Among attackers that evade
in at least 5\% of missions, the largest undetected contribution lies between
\nDevSingleMin{} and \nDevSingleMax\,m \emph{whichever detector is deployed}.
Rarer evasions go further: up to \nGatedUnfMax\,m against \nGatedUnfStacks{}
(\nGatedUnfFam). The pooled supremum against the full stack is \nAnchorMaxFive\,m
(\cref{fig:anchor}), and the stealth budget is $\dadv=\nDadvFive$\,m. Within that
range the trusted fix, not the detector design, decides how far an undetected
attacker gets. The median KS distance between a detector's statistic under an
evading attack and under benign conditions is \nKsStealthyMed{}, against
\nKsCaughtMed{} for detected ones. This is the empirical face of
\cref{lem:lecam}: an attack that evades is, to the detector, almost
indistinguishable from the benign population.

\paragraph{Without the gate, each detector admits a different disguise.}
The bottom of \cref{tab:disguise} repeats the cross-evaluation with the anchor
still fused into the filter but its alarm removed, isolating what the detectors
themselves discriminate. Evasion stays at or above \nNgEvMin{}, but the
undetected contribution now depends sharply on the pairing of disguise and
detector.
\begin{itemize}
\item The fault mimic drives \nNgDOneFaultMimic{}\,m past D1, D3 and D4. Its slow,
  bias-shaped curvature has no jumps for D1, matches the bias drift that D4
  tolerates by design, and looks like an unmodelled disturbance to D3.
\item The GNSS mimic and the weather mimic exploit the two detectors that do
  not cross-check the IMU, D1 and D3 (up to \nNgDOneGnssMimic{} and
  \nNgDOneWeatherMimic\,m). Among frequently evading attackers, D4 caps them at
  \nNgDFourWeatherMimic\,m or below, because their acceleration content
  contradicts the accelerometer; a rare GNSS-mimic evasion still reached
  \nNgDFourGnssUnf\,m.
\item Rare step offsets that happen to land in a benign noise trough evade D1,
  D3 and D4 in about 3\% of missions and reach \nNgStepUnfMax\,m.
\item Only the $\chi^2$ gate D2, whose filter models the accelerometer bias,
  caps every family, at \nNgCapDTwo\,m (fault mimic), including rare evasions
  (\nNgDTwoUnfMax\,m); the full stack, which contains D2, is capped at
  \nNgCapALL\,m.
\end{itemize}
This is the pattern adversarial equivalence predicts. The disguise that wins
is the one the detector models as benign: bias drift for an invariant that
tolerates bias, and drift or wind-shaped motion for detectors without an
inertial cross-check.

\begin{table}[t]
\centering\footnotesize
\setlength{\tabcolsep}{3pt}
\caption{Disguise matrix at 5\% FPR, cross-evaluated over all
\nPooledAttackers{} optimised attackers on \nEvalMissions{} held-out missions
per attacker. \emph{Top:} anchor gate armed; best evasion rate (\%) / largest
undetected attack contribution (m). \emph{Bottom:} anchor fused but its gate
removed; largest undetected contribution (m). Contributions are taken over
attackers that evade in at least 5\% (and at least 10) of missions; rarer
evasions are reported in the text. Benign missions raise no alarm in about
\nBenignNoAlarm{}. Darker cells are evaded more often (top) or admit more
deviation (bottom).}
\label{tab:disguise}
\begin{tabular}{@{}lccccc@{}}
\toprule
& \textbf{step} & \textbf{ramp} & \textbf{fault} & \textbf{GNSS} & \textbf{weather} \\
& & & \textbf{mimic} & \textbf{mimic} & \textbf{mimic} \\
\midrule
\multicolumn{6}{@{}l}{\emph{anchor gate armed (C2$'$)}}\\
\tabinput{tab_disguise_body.tex}
\midrule
\multicolumn{6}{@{}l}{\emph{anchor fused, gate removed}}\\
\tabinput{tab_disguise_nogate_body.tex}
\bottomrule
\end{tabular}
\end{table}

\paragraph{The terminal window.}
\Cref{fig:anchor} sweeps the re-anchoring interval against the full stack,
pooling all attackers on common missions. With every interval closed, the
largest undetected contribution grows from \nAnchorMaxFive\,m at $T_a=5$\,s to
\nAnchorMaxTwenty\,m at 20\,s (\nAnchorMultTwenty$\times$). The attacker must
return near zero before each trusted fix, so the budget is set by how far it can
travel and come back within one interval. When the segment before mission end
is left unchecked, no return leg is needed, and the contribution reaches
\nOpenMaxTwenty\,m, \nOpenRatioTwenty$\times$ the closed value. The terminal
window is where an equivalence attacker spends its budget.

\figAnchor

\subsection{RQ3: where the looseness lives}
\Cref{fig:gap} decomposes the tail of the benign budget; the ratios below are
at $s=\nGapS$\,m, and their ordering changes in the far tail, where few joint
exceedances remain.
\begin{enumerate}
\item For two sources, the union and RA bounds agree to within
  \nUnionRaDtwoMaxErr{}, as Makarov's theorem requires. For three, the RA
  tightens the union bound by \nRaGain{} (\nGapUnion{} to \nGapRa{}). Under
  unknown dependence, the union bound is nearly the right quantity.
\item The worst coupling exceeds independent sums at the same marginals by
  \nDepRatio$\times$ (\emph{dependence gap}). Evaluating the wind at the upper
  envelope rather than at the conditions actually flown costs a further
  \nEnvRatio$\times$ (\emph{envelope gap}).
\item Independent sums at matched wind agree with the sum of per-source
  suprema measured on the same missions to within \nNatRatio$\times$,
  consistent with the benign sources being independent, as they are by
  construction in the simulator.
\item The sum of per-source suprema exceeds the supremum of the joint deviation
  by \nTimeRatio$\times$ (\emph{temporal gap}; \nGapJointK{} joint exceedances in
  \nJointMissions{} missions). Sources seldom peak at the same instant, but a
  budget built from per-source suprema must assume they do.
\end{enumerate}
The dependence gap can be closed only by \emph{knowing} the coupling, which an
attacker who controls jamming denies the defender. The envelope gap shrinks
with better wind knowledge. The temporal gap, the largest of the three at
this $s$, can be
closed without either, by bounding the supremum of the summed trajectory
directly, for example with a reachable-set computation per certification
window. Condition (C5), which licenses adding per-source suprema, held in
\nSubaddSources{} of joint missions.

\figGap

\subsection{RQ2 and RQ4: the bound under evasion}
Pooling all attackers on calibration missions gives a stealth budget of
$\dadv=\nDadvFive$\,m, attained by the \nDadvArgFive. The budget is identical at
1\%, 5\% and 10\% FPR. It is a single pooled maximum over common missions, so this
says that no attacker in the pool gains from looser thresholds in that range,
not that thresholds never matter.

\Cref{tab:soundness} reports, for each operating point, the bounds on
$\Pr[\text{fail}\wedge U]$ and the largest observed frequency of that event over
five held-out attackers and three wind conditions. Across \nSoundPoints{}
distinct operating points and \nUndetectedTotal{} undetected-attack missions,
\textbf{no point shows a Holm-significant violation} (\nSoundViol{} of
\nSoundPoints). The scaled test of \cref{sec:scaled} overturns this. Wherever the bound is informative ($\pbar<1$), the observed
frequency of undetected failure was zero. The bound is therefore conservative,
and the zero violations by themselves carry little evidence: the negative
controls below are what show that the test can fail.

Three further checks belong with the result.
\begin{itemize}
\item \emph{The empirical budget is not a theorem.} The held-out attacker
  exceeded $\dadv$ in \nSupExceed{} of \nSupChecks{} (family $\times$ FPR $\times$
  wind) cells, reaching \nHeldoutAdvMax\,m with the \nHeldoutAdvFam{} at
  \nHeldoutAdvFpr{} FPR (\nHeldoutAdvRatio$\times$). The bound absorbed it, but
  this is why we strengthened the optimisers and derived an analytic budget
  (\cref{sec:scaled,sec:revised}). In \nZeroUndet{} cells, the
  held-out optimiser found no evading attack at all.
\item \emph{\Cref{lem:subadd} holds where (C1) holds.} Attack/benign
  superposition held in at least \nSuperposMin{} of missions at every point. A
  control run with the actuator limit removed (5\% FPR, envelope wind, all five
  families) restored it in \nSatUnsatMin{} of missions (largest excess
  \nSatUnsatExcess\,m). We attribute the larger excesses at other points, up to
  \nSuperposExcessMax\,m, to the same mechanism, saturation, which (C1)
  excludes; they were not re-run.
\item \emph{Independence is not certified.} The independence bound
  $p_\perp$ was not violated either, because the simulated sources are
  independent. It would claim a 0.9 floor at $m=\nMNinetyIndep$\,m, against
  \nMNinety\,m for the certified bound (\cref{fig:floor}); nothing guarantees
  that independence when an attacker controls jamming.
\end{itemize}

\begin{table}[t]
\centering\footnotesize
\setlength{\tabcolsep}{2.2pt}
\caption{Soundness of \cref{thm:floor} on held-out attackers. Bounds on
$\Pr[\text{fail}\wedge U]$: union $p_U$, RA $\pbar$ (certified) and independence
$p_\perp$ (not certified). Obs.\ = largest observed frequency of
$\{\text{fail}\wedge U\}$ over five families and three wind conditions;
MCR$_U$ = smallest completion rate among undetected missions; viol.\ =
Holm-significant violations of $\pbar$ and of $p_\perp$. $\dadv=\nDadvFive$\,m.
Rows with $\theta_{\mathrm{net}}=0.25$\,s shift the margin analytically by
$e^{LT_c}\dnet=\nNetTubeQuarter$\,m; the network term itself is not simulated.
$p_U$ is a proven upper bound; the RA column is a numerical approximation of
the sharp $\pbar$, reported as $\min(\mathrm{RA},p_U)$.}
\label{tab:soundness}
\begin{tabular}{@{}cccccccccc@{}}
\toprule
\textbf{FPR} & $\theta_{\mathrm{net}}$ & $m$ & $p_U$ & $\pbar$ & $p_\perp$ & \textbf{obs.} & \textbf{MCR$_U$} & \multicolumn{2}{c}{\textbf{viol.}} \\
\midrule
\tabinput{tab_soundness_body.tex}
\bottomrule
\end{tabular}
\end{table}

\figFloor

\paragraph{Negative controls: does the test have power?}
A soundness test that cannot fail is uninformative, so we broke two conditions
deliberately and re-optimised the attacker under each (\cref{tab:negative}).
\begin{itemize}
\item \emph{(C2$'$) broken, no re-anchoring.} Undetected deviation reaches
  \nNegCtwoMaxDev\,m, and the test flags \nNegCtwoViol{} of \nNegCtwoRows{}
  points (\nNegCtwoViolFam{}), every one where the bound is informative.
\item \emph{(C4) broken, severe turbulence (45--60\,kt).} The fault mimic
  reaches \nNegCfourMaxDev\,m undetected, and the test flags the bound at
  $m=20$\,m, where $\pbar=0$ because the empirical marginals have finite
  support. At $m=10$\,m the observed frequency stays below $\pbar$, so this
  break is detected only where the bound is tightest. The weather mimic was
  detected in every mission, so there was nothing to test.
\end{itemize}
The test thus detects unsoundness when a stated condition fails and enough
undetected missions remain to measure it, which is the behaviour the audit
loop of \cref{fig:framework} relies on.

\begin{table}[t]
\centering\footnotesize
\setlength{\tabcolsep}{2.4pt}
\caption{Negative controls at 5\% FPR. Undet.\ = undetected missions out of
\nEvalMissions; fails = failures among them; viol.\ = Holm-significant violation
of $\pbar$; dev.\ = largest undetected deviation (m).}
\label{tab:negative}
\begin{tabular}{@{}llcccccc@{}}
\toprule
\textbf{Control} & \textbf{disguise} & $m$ & \textbf{undet.} & \textbf{fails} & $\pbar$ & \textbf{viol.} & \textbf{dev.} \\
\midrule
\tabinput{tab_negative_body.tex}
\bottomrule
\end{tabular}
\end{table}

% ------------------------------------------------------------- Scaled test (R1) and revision
\subsection{RQ5: the test at scale breaks the pilot instantiation}
\label{sec:scaled}
The pilot's zero violations rested on five CEM attackers per cell and a
conservative bound. We therefore strengthened the attacker and the test before
drawing conclusions (runbook R1). \nROneAttackers{} further attackers were
optimised against the full stack: CMA-ES (population \nROneCmaPop,
\nROneCmaGens{} generations, one restart)~\cite{hansen2016cma}, Bayesian
optimisation with a tree-structured Parzen estimator (\nROneTrials{}
trials), and the same TPE with a tail objective (the 90th percentile of the
undetected deviation instead of its mean). Every candidate was scored on
\nROneK{} missions with common random numbers, and a sixth \emph{mixed} family
switches between any two disguises mid-mission. The scaled soundness test
re-optimised \nROneHeldSeeds{} independent CEM attackers per (family, FPR) cell
and evaluated each on \nROneEvalN{} fresh missions under three wind
conditions: \nROneTests{} tests over \nROneUndet{} undetected-attack missions.

\paragraph{The stealth budget grew by 25\%.}
The pooled budget rose from \nDadvFive{} to \nROneDadvFive\,m at 1\% and 5\% FPR
(\nROneDadvRatioFive$\times$) and to \nROneDadvTen\,m at 10\%, attained by the
\nROneDadvArg. Only the tail objective found more than the pilot's CEM: the
mean-objective CMA-ES and TPE peaked at \nROneSupCma{} and \nROneSupTpe\,m. An
attacker that optimises its typical deviation is not the attacker that sets an
empirical supremum. Even the larger budget was exceeded by a held-out attacker
in \nROneSupExceed{} of \nROneCells{} cells, by the mixed family up to
\nROneAdvMax\,m. An empirical supremum is a lower bound on $\dadv$, and more
search keeps raising it.

\paragraph{Five Holm-significant violations.}
With $\dadv=\nROneDadvFive$\,m, the test flagged \nROneViol{} of \nROneTests{}
points, all at $m=\nROneViolM$\,m, where the bound claimed $\pbar=0$. We
re-simulated each flagged cell with its seeds
(\code{r1\_violation\_diagnosis.csv}). The violations come from
\nDiagMissions{} missions, and the attacker caused none of them. Each fails
under its benign inputs alone (benign-only deviation \nDiagBenignMin--\nDiagBenignMax\,m),
the attack adds at most \nDiagAttackMax\,m, and \nDiagAllSat{} three saturate the
actuator with or without the attack. Two conditions of \cref{thm:floor} failed,
and neither is about disguise.
\begin{itemize}
\item \emph{(C3) for the marginals.} The empirical marginal of each source has
  finite support. Beyond the largest of \nCalMissions{} calibration samples,
  $\pbar$ is exactly zero, and the bound claims certainty that no sample can
  justify. One benign mission in 5{,}000 fails at $m=20$\,m.
\item \emph{(C1), which nothing enforced.} Superposition held in at least
  \nROneSuperposMin{} of missions, and every violating mission saturated. The
  pilot had stated (C1) as an assumption, but no component of the stack checked
  it at run time.
\end{itemize}
This is a negative result about the pilot instantiation, not about
\cref{thm:floor}. The theorem's conclusion requires its conditions, and the
scaled test found two that the instantiation did not secure. We report the
violation as found. The revision below fixes the instantiation before looking at
new data and does not tune the old run until it passes.

\subsection{A revised instantiation, judged on fresh data}
\label{sec:revised}
We changed three things, fixed them before generating any new data, and
evaluated them on a new seed block. A smoke test of the code used a separate
block, which we did not reuse.
\begin{enumerate}
\item \emph{A saturation alarm $D_{\mathrm{sat}}$.} The autopilot raises an alarm
  at any epoch at which the commanded acceleration reaches its limit. This
  turns (C1) from an assumption into a run-time check: on the event $U$, the
  loop never saturated, so \cref{lem:subadd} applies exactly. The alarm costs
  benign missions \nRevSatAlone{} in false alarms (5\% operating point, envelope
  wind). The full stack's false-alarm rate rises from \nRevFprStack{} to
  \nRevFprSat.
\item \emph{A finite-sample p-box for (C3).} Each marginal tail is replaced by
  its Clopper--Pearson upper bound, taken simultaneously over the split grid
  (Bonferroni over $3\times\nRevGrid$ grid points, $\alpha=0.05$), and
  aggregated by the union bound. With probability at least 0.95 over the
  calibration samples, $\pbar$ is then a valid upper bound everywhere on the
  grid. It never reaches zero, and with \nCalMissions{} samples per source it
  levels off at $1-\nRFiveFloorCpTwenty$.
\item \emph{Two stealth budgets, tested separately.} The first is the R1
  calibration supremum, $\dadv=\nROneDadvFive$\,m. It was \emph{not} widened to
  the \nROneAdvMax\,m that the R1 held-out attacker reached, because a constant
  widened after looking at held-out data would make the new test circular. The
  second is the analytic bound below, which holds for every attack.
\end{enumerate}

\paragraph{An analytic stealth budget (R8).}
With noise removed, the linear loop maps the attack offsets $d_j$ at the GNSS
epochs to the attack-only deviation linearly. The worst undetected deviation
is therefore a linear programme over $d$, solved once for every time step.
Its constraints are the detector conditions that are convex in $d$ and
\emph{implied} by silence, with the benign contribution bounded by its
$1-\epsilon/3$ calibration quantile ($\epsilon=0.01$).
\begin{itemize}
\item D1: a bounded jump between consecutive fixes.
\item The anchor gate: a bounded innovation at each trusted fix.
\item $D_{\mathrm{sat}}$: a bounded attack-only command.
\end{itemize}
D2--D4 alarm on runs of exceedances or on quadratic statistics, which are not
convex in $d$, so we drop them. Dropping a detector only enlarges the feasible
set, so the bound stays valid for every stack that contains the kept ones. We
take the worst actuator effectiveness in $[0.7,1]$, and the onset of the attack
is free. Three findings follow (\code{r8\_analytic\_bound.csv}).
\begin{itemize}
\item \emph{The anchor gate alone bounds nothing.} The trusted fix constrains
  the estimate error at anchor times, not the deviation. An attacker that
  carries the estimate along with the vehicle passes every anchor.
\item \emph{A jump detector adds little.} With D1 the budget is finite but
  \nREightDOneFive\,m at $T_a=5$\,s, because D1's threshold must tolerate
  the jumps of GNSS jamming noise.
\item \emph{The saturation alarm is what makes it physical.} With
  $D_{\mathrm{sat}}$ the budget falls to \nREightSatFive\,m. That is still
  \nREightSatRatio$\times$ the calibrated $\dadv$ and \nREightSatRatioAdv$\times$
  the largest deviation any optimised attacker reached (\nROneAdvMax\,m).
\end{itemize}
The gap is the measure of what the non-convex detectors D2--D4 contribute and
what no one has yet proved. Using the provable budget, the floor certifies
nothing below a margin of about \nREightSatFive\,m. Using the empirical budget, it
certifies more, but rests on a supremum that more search has raised before.

\paragraph{Result on fresh seeds.}
We re-optimised \nRevAttackers{} held-out attackers against the revised stack:
three CEM seeds per (family, FPR) cell plus one tail-objective TPE. Each was
evaluated on \nROneEvalN{} missions under three wind conditions, from a seed
block never used before (\nRevUndet{} undetected-attack missions).
\begin{itemize}
\item \emph{No violation under either budget.} With the R1 budget, the test found
  \nRevEmpViol{} Holm-significant violations in \nRevEmpTests{} tests; with the
  LP budget, \nRevLpViol{} in \nRevLpTests. Superposition was exact on every
  undetected mission (largest excess \nRevSupUnsat\,m). Missions like the benign-tail failures that broke the pilot instantiation
  saturate, so they now raise the alarm and leave the event $U$.
\item \emph{The empirical budget was exceeded again.} In \nRevEmpExceed{} of
  \nRevEmpCells{} cells a held-out attacker went beyond $\dadv=\nROneDadvFive$\,m,
  up to \nRevAdvMax\,m (\nRevAdvMaxFam, \nRevAdvMaxOpt). Benign slack absorbed it
  and no violation followed, but the stealth budget we can use is still not a
  sound constant. The LP budget was never exceeded. It can be exceeded only with probability
  $\epsilon$, or on missions with a GNSS outage, which its proof does not cover.
\item \emph{What the bound certifies.} With the R1 budget, the floor at 5\% FPR
  is \nRevEmpFloorTen{} at $m=10$\,m and \nRevEmpFloorTwenty{} at $m=20$\,m.
  With the LP budget it is \nRevLpFloorTen{} at both, so it certifies nothing
  in the corridors we test.
\item \emph{Limited power.} Of the \nRevEmpInformative{} tests at which the bound
  is informative, \nRevEmpFailInf{} saw any undetected failure. The revision is
  therefore \emph{not refuted}, not \emph{confirmed tight}. The negative
  controls remain the evidence that the test can fail, and the scaled run of
  \cref{sec:scaled} shows that it does when a condition is missing.
\end{itemize}


\subsection{RQ6: how much does a measured GNSS marginal cost?}
\label{sec:r5}
The pilot's GNSS marginal is simulated (99th percentile \nRFivePnnSim\,m). The
parent artifact contains measured navigation error during real GPS jamming:
\nRFiveSamples{} one-second samples over two episodes, a noise episode and a
loss of fix, whose error envelope grows with staleness to \nRFiveEnvTen\,m at
10\,s and \nRFiveEnvMax\,m at 15\,s (runbook R5; spoofing samples are excluded,
since spoofing belongs to $\dadv$). We replaced each simulated degradation episode
by the larger of its simulated deviation and the measured envelope at the
episode's staleness, with the simulator's episode probabilities, and in a stress
variant set every episode to the measured 45\,s duration. The construction is
conservative in one respect, since it never makes a mission's GNSS deviation
smaller. It also assumes that the measured navigation error becomes cross-track
deviation one for one, which holds when the controller tracks its estimate
and is not itself a measurement.

At $m=10$\,m and 5\% FPR, with $\dadv=\nROneDadvFive$\,m, the floor estimated from
the empirical marginals falls from \nRFiveFloorSim{} (simulated GNSS) to
\nRFiveFloorMeas{} (measured envelope) and \nRFiveFloorLong{} (45\,s episodes).
With the finite-sample band of \cref{sec:revised}, the values are
\nRFiveFloorCpSim, \nRFiveFloorCpMeas{} and \nRFiveFloorCpLong. The bound is only
as good as its GNSS marginal. A simulated marginal overstated the certified
floor by a factor of two to three at this margin. The jamming-induced tail,
not the attack, is what the corridor must absorb. At $m=20$\,m every variant
certifies \nRFiveFloorCpTwenty, the ceiling that \nCalMissions{} calibration
samples per source permit.


\subsection{Summary}
\textbf{RQ1:} every disguise family evades every detector; what differs is the
damage. With the anchor gate armed, frequent evasions stay within
\nGatedCap\,m and the pooled supremum is \nAnchorMaxFive\,m, whatever the
detector. Without the gate, the damage depends on whether the detector models
the disguise, from \nNgCapDTwo\,m (D2) to \nNgMax\,m (D1, D3, D4 against the
fault mimic).
\textbf{RQ2 and RQ4:} in the pilot, the bound held at all \nSoundPoints{}
held-out points, and the negative controls showed that the test can fail.
\textbf{RQ3:} under unknown dependence the union bound is nearly sharp; the
slack lies in the dependence, envelope and, above all, temporal gaps.
\textbf{RQ5:} at scale, the pilot instantiation was unsound
(\nROneViol{} violations). The cause was not disguise but two conditions it
assumed without securing: a finite-sample marginal and an unchecked actuator
limit. Once both were secured, the bound survived \nRevAttackers{} fresh
attackers with \nRevEmpViol{} violations, although the empirical $\dadv$ was
exceeded again. The provable budget is \nREightSatRatio$\times$ larger.
\textbf{RQ6:} measured jamming data halves the certified floor at
$m=10$\,m: a simulated GNSS marginal overstates the floor by a factor of two.
